% =============================================================
% Abstract, key words and abbreviations (from PaperDINOcow.pdf)
% =============================================================
\section*{Abstract}

Body condition scoring (BCS) is used to assess body fat reserves in dairy cattle.
However, manual scoring \mP{is time-consuming and costly,} requires repetitive observations and sometimes varies between evaluators, making continuous monitoring difficult.
This study presents DINOcow, a computer vision approach for continuous BCS estimation using two-dimensional images captured in a robotic milking system.
A dataset of 670 rear dorsal images from Holstein dairy cows was used to develop and evaluate the model.
This model combines two self-supervised vision transformers.
First, a DINOv2-B/14 model with a feature pyramid decoder segments six anatomical regions associated with body condition assessment.
Then, a DINOv3-L/16 model extracts visual features from the complete dorsal image.
The segmentation masks are used to aggregate these features by anatomical region, and five regional descriptors are combined to estimate BCS through ridge regression.
The performance of the model was evaluated using leave-one-cow-out (LOCO) and repeated five-fold cross-validation, keeping images from the same cow within the same data partition.
Under LOCO evaluation, the model achieved \mE{an MAE} of 0.415 BCS points across images and an average MAE of 0.405 per cow, with an $R^2$ of 0.72 and a Pearson correlation coefficient of 0.85.
The use of anatomical regions improved BCS estimation compared with descriptors obtained from the whole image, the cow silhouette, or a single rear crop.
These results show that combining anatomical segmentation with frozen visual representation can support BCS estimation using conventional two-dimensional images, while reducing the need to train large models with extensive BCS-labeled datasets.

\medskip
\noindent\textbf{Key words:} body condition score, vision transformers, feature extraction, dairy cattle, anatomical segmentation

\noindent\textbf{Nonstandard abbreviations:}\mJ{AUC = area under the ROC curve; IoU = intersection over union; LOCO = leave-one-cow-out; LoRA = low-rank adaptation; MAE = mean absolute error; mIoU = mean IoU; PIR = passive infrared; QWK = quadratically weighted kappa; RMSE = root mean squared error; ROC = receiver operating characteristic; ViT = vision transformer}

